% !TeX root = ../main.tex
\section{Решение задачи 3-4}

\vspace{1em}
\textbf{1. Исходные данные}

Из условия задачи известны следующие параметры:
\begin{itemize}
    \item Девиация частоты: $\Delta f = 125 \text{ кГц} = 1.25 \cdot 10^5 \text{ Гц}$
    \item Частота модуляции: $F_{\text{м}} = 25 \text{ кГц} = 2.5 \cdot 10^4 \text{ Гц}$
    \item Входное сопротивление БВЧ: $R_{\text{вх}} = 50 \text{ Ом}$
    \item Амплитуда немодулированного сигнала: $U_m = 10 \text{ мкВ} = 10 \cdot 10^{-6} \text{ В}$
    \item Спектральная плотность мощности шума: $N_0 = 4 \cdot 10^{-19} \text{ Вт/Гц}$
\end{itemize}

\vspace{1em}
\textbf{2. Расчет параметров сигнала и шума на входе}

Индекс частотной модуляции рассчитывается как отношение девиации к частоте модуляции:
\begin{align*}
    \beta_{\text{чм}} &= \frac{\Delta f}{F_{\text{м}}} = \frac{125}{25} = 5
\end{align*}

Мощность полезного сигнала на входе приемника:
\begin{align*}
    P_c &= \frac{U_m^2}{2 R_{\text{вх}}} = \frac{(10 \cdot 10^{-6})^2}{2 \cdot 50} = \frac{100 \cdot 10^{-12}}{100} = 10^{-12} \text{ Вт}
\end{align*}

АЧХ БВЧ и БНЧ прямоугольные и согласованы по полосе с соответствующими сигналами. Полоса пропускания высокочастотного тракта (БВЧ) для ЧМ сигнала приведена к удвоенной девиации частоты:
\begin{align*}
    \Pi_{\text{БВЧ}} &= 2 \Delta f = 2 \cdot 125 \text{ кГц} = 250 \text{ кГц} = 2.5 \cdot 10^5 \text{ Гц}
\end{align*}

Мощность шума в полосе БВЧ на входе детектора:
\begin{align*}
    P_{\text{ш.вх}} &= N_0 \cdot \Pi_{\text{БВЧ}} = 4 \cdot 10^{-19} \cdot 2.5 \cdot 10^5 = 10 \cdot 10^{-14} = 10^{-13} \text{ Вт}
\end{align*}

Отношение сигнал-шум на входе (в разах и децибелах):
\begin{align*}
    q_{\text{вх}} &= \frac{P_c}{P_{\text{ш.вх}}} = \frac{10^{-12}}{10^{-13}} = 10 \\
    q_{\text{вх, дБ}} &= 10 \log_{10}(10) = 10 \text{ дБ}
\end{align*}

\vspace{1em}
\textbf{3. Расчет отношения сигнал-шум на выходе}

Согласно формуле идеальной модели (надпороговый режим), представленной на графике (красная пунктирная линия), отношение сигнал-шум на выходе ЧМ-приемника рассчитывается как:
\begin{align*}
    q_{\text{вых}} &= 3 \beta_{\text{чм}}^3 \cdot q_{\text{вх}}
\end{align*}

Подставляем найденные значения:
\begin{align*}
    q_{\text{вых}} &= 3 \cdot 5^3 \cdot 10 = 3 \cdot 125 \cdot 10 = 3750 \text{ (в разах)}
\end{align*}

Переводим результат в децибелы для получения окончательного ответа:
\begin{align*}
    q_{\text{вых, дБ}} &= 10 \log_{10}(3750) \approx 35.74 \text{ дБ}
\end{align*}

\textbf{Примечание к графику:}
Как видно из представленного графика, при $q_{\text{вх}} = 10 \text{ дБ}$ и $\beta_{\text{чм}} = 5$ рабочая точка находится на границе наступления порогового эффекта. Идеализированная расчетная модель дает $35.74 \text{ дБ}$ (пунктирная линия), в то время как реальная характеристика (сплошная синяя линия модели) начинает незначительно отклоняться вниз из-за воздействия шумов.

\textbf{Ответ:} Рассчитанное отношение сигнал-шум на выходе приемника составляет $\approx 35.74 \text{ дБ}$.